\documentclass[12pt]{article}
\usepackage{graphicx}
\usepackage{hyperref}  % For URLs
\usepackage[a4paper, total={6in, 8in}]{geometry}  % Adjust page margins
\usepackage{setspace}  % For line spacing
\usepackage{parskip}   % To remove extra space between paragraphs
\usepackage{tabularx}  % To create flexible tables
\usepackage{multirow}  % For multi-row cells in tables

% Adjust line spacing
\linespread{1.1}

\title{Natural Vision Restoration: Experimental Verification, Astigmatism Reduction, and the ``Just Barely Clear'' Principle}
\author{Li Zhiwei}
\date{August 2026}

\begin{document}

\maketitle

\tableofcontents

\section*{Abstract}
This paper presents a unified account of the natural vision restoration method for reversing myopia, combining three strands of research based on the author's personal experiments: (1) a one-year experimental verification of the method, which produced significant reductions in both myopia and astigmatism; (2) an analysis of the role of astigmatism during eyeball deformation and recovery, showing that astigmatism reduction is an integral part of the overall reversal process and can be measured as an equivalent myopic reduction; and (3) the discovery of the ``just barely clear'' principle, which holds that the key to continued reversal lies not in a fixed degree of prescription reduction but in consistently maintaining a state where objects are sufficiently clear without strain. Together, these results support the hypothesis that myopia and astigmatism are partially reversible conditions, driven by the gradual reshaping of the eyeball.

\section{Introduction}
Myopia, or nearsightedness, is a common refractive error in the human eye. Traditionally, corrective lenses are prescribed to bring distant objects into focus, but the concept of myopia reversal---restoring the eye's natural shape to reduce dependence on corrective lenses---is gaining interest. Astigmatism, another form of refractive error caused by the uneven curvature of the cornea or lens, is often present alongside myopia. While myopia has been well-studied, the role of astigmatism in myopia correction and reversal remains less explored, and the practical principles that make reversal sustainable over time are not widely known.

This paper combines the author's earlier investigations into a single comprehensive study of natural vision restoration. It presents the results of an experiment, begun in March 2022, that verifies the effectiveness of the natural vision restoration method inspired by the works of Todd Becker (2014) and Yin Wang (2022). It then examines the relationship between reductions in myopia and astigmatism during the process of eyeball deformation and recovery, and finally introduces the ``just barely clear'' principle---a novel, personalized approach to selecting corrective lenses that was found to be crucial for continued progress.

\section{Background}
\label{sec:background}
Myopia is a prevalent condition that has been on the rise globally. Its causes are multifaceted, involving both genetic and environmental factors. Notably, modern lifestyles---characterized by prolonged near work (e.g., reading, screen time)---have been implicated in the increase of myopia cases.

According to Todd Becker's method, first introduced in 2014, myopia can be reversed by gradually reducing the strength of corrective lenses, thereby encouraging the eye muscles to regain their natural shape. An article titled ``Myopia: A Modern Yet Reversible Disease,'' presented at the 2014 Ancestral Health Symposium, highlighted that myopia is a reversible condition. The theory posits that myopia progresses in two stages: (1) near work induces lens spasm (pseudo-myopia), and (2) using minus lenses temporarily corrects the vision, leading to eye elongation and further myopia.

Yin Wang's more recent work (2022) built on Becker's foundation, offering insights into the biological mechanisms underlying this approach. He proposed a mechanism by which reducing the strength of corrective lenses promotes the relaxation and reshaping of eye muscles, ultimately leading to improved vision. A key passage from Wang's work, which proved to be a turning point in the author's understanding, describes the mechanism of optical axis shortening:

\begin{quote}
    When looking at distant objects, the crystalline lens needs to relax and reduce refraction. If the image still falls in front of the retina after the lens is completely relaxed, objects appear blurry. If you continue to look at the object in this state, the external eye muscles will begin to gently compress the eyeball, making the optical axis slightly shorter (perhaps only 0.x millimeters), and the blurry object will gradually become somewhat clearer. If you frequently maintain this slightly blurry state, the repeated small compressions of the optical axis will lead to permanent shortening, and myopia will reverse.
\end{quote}

This passage clarified the underlying mechanism by which the optical axis may shorten through gradual, repeated compression, and it motivated the practical experiments described in this paper.

\section{Methodology}
The author has tracked their own myopia and astigmatism progression since March 2022, using a corrective eyewear method designed to reduce the curvature of the eyeball over time. The method involves wearing eyeglasses with a degree of myopia lower than the actual prescription and monitoring vision clarity at typical viewing distances, such as for mobile phone use and computer work.

The methodology includes:
\begin{itemize}
    \item Wearing corrective lenses with varying degrees of reduction (ranging from 150 to 200 degrees).
    \item Monitoring vision clarity during daily activities.
    \item Recording eye examination data at regular intervals.
    \item Noting improvements or stagnation in myopia and astigmatism.
    \item Analyzing the experience and correlating it with the relevant theory in myopia reversal.
\end{itemize}

The data from eye examinations between March 2022 and April 2023 are used to explore the relationship between reductions in myopia and astigmatism.

\subsection{Initial Report}
The following table presents the myopia and astigmatism measurements taken on March 5, 2022:

\begin{table}[h!]
\centering
\begin{tabularx}{\textwidth}{|X|X|X|}
\hline
\textbf{Date} & \textbf{Left Eye (°)} & \textbf{Right Eye (°)} \\
\hline
2022.03.05 & Myopia: 350, Astigmatism: 225 & Myopia: 575, Astigmatism: 175 \\
\hline
\end{tabularx}
\end{table}

\subsection{Subsequent Reports}
The following table shows the myopia and astigmatism measurements on November 13, 2022, and April 20, 2023, respectively:

\begin{table}[h!]
\centering
\begin{tabularx}{\textwidth}{|X|X|X|}
\hline
\textbf{Date} & \textbf{Left Eye (°)} & \textbf{Right Eye (°)} \\
\hline
2022.11.13 & Myopia: 325, Astigmatism: 200 & Myopia: 550, Astigmatism: 175 \\
\hline
2023.04.20 & Myopia: 300, Astigmatism: 125 & Myopia: 500, Astigmatism: 125 \\
\hline
\end{tabularx}
\end{table}

\section{Results}

\subsection{Myopia and Astigmatism Reduction}
Over the course of the experiment, significant improvements were observed:
\begin{itemize}
    \item My left eye's myopia reduced from 350 degrees to 300 degrees (a decrease of 50 degrees).
    \item My right eye's myopia reduced from 575 degrees to 500 degrees (a decrease of 75 degrees).
    \item Both eyes showed improvement in astigmatism, with the left eye's astigmatism reducing from 225 degrees to 125 degrees, and the right eye's astigmatism reducing from 175 degrees to 125 degrees.
\end{itemize}

The following table summarizes the eye examination results for the author's left and right eyes over the course of the year:

\begin{table}[h!]
\centering
\begin{tabularx}{\textwidth}{|X|X|X|X|X|}
\hline
\textbf{Time} & \textbf{Myopia (Left)} & \textbf{Astigmatism (Left)} & \textbf{Myopia (Right)} & \textbf{Astigmatism (Right)} \\
\hline
2022.03 & 350 & 225 & 575 & 175 \\
2023.04 & 300 & 125 & 500 & 125 \\
\hline
\textbf{Reduction} & \textbf{50} & \textbf{100} & \textbf{75} & \textbf{50} \\
\hline
\end{tabularx}
\end{table}

These results were unexpected and highly encouraging, confirming that the natural vision restoration method had a noticeable effect on both conditions.

\subsection{Eyewear Adjustment}
To better understand the results, I carefully tracked my eyewear changes throughout the experiment. In November 2022, I switched to a new pair of glasses with 150 degrees less than my prescribed myopia, following the recommendations from Todd Becker's method.

Here is the comparison of my eyesight and eyeglasses prescription on November 13, 2022:

\begin{table}[h!]
\centering
\begin{tabularx}{\textwidth}{|X|X|X|}
\hline
\textbf{Item} & \textbf{Left Eye (°)} & \textbf{Right Eye (°)} \\
\hline
My eyesight & Myopia: 325, Astigmatism: 200 & Myopia: 550, Astigmatism: 175 \\
Glasses I was wearing & Myopia: 225, Astigmatism: 200 & Myopia: 450, Astigmatism: 175 \\
\hline
\end{tabularx}
\end{table}

The reduction in myopia and astigmatism led to further improvements over the next six months. On April 20, 2023, my updated prescription was as follows:

\begin{table}[h!]
\centering
\begin{tabularx}{\textwidth}{|X|X|X|}
\hline
\textbf{Item} & \textbf{Left Eye (°)} & \textbf{Right Eye (°)} \\
\hline
My eyesight & Myopia: 300, Astigmatism: 125 & Myopia: 500, Astigmatism: 125 \\
Glasses I was wearing & Myopia: 175, Astigmatism: 200 & Myopia: 400, Astigmatism: 175 \\
\hline
\end{tabularx}
\end{table}

\section{Astigmatism Reduction During Myopia Reversal}
The effect of astigmatism reduction is quite significant during the process of myopia reversal. Notably, the reduction in astigmatism should be considered alongside the reduction in myopia. Astigmatism, like myopia, is a measure of the eye's refractive power, and its reduction is not merely a secondary condition but an integral part of the overall process of eyeball deformation and recovery. It is suggested that the degrees of astigmatism reduction should be added to the degrees of myopia reduction because the curvature of the eye that leads to astigmatism is also involved in the development of myopia. Dividing the astigmatism degrees by two provides a reasonable approximation of its equivalent myopic reduction.

From the data above, we can calculate the total reduction of both myopia and astigmatism in both eyes:

\begin{table}[h!]
\centering
\begin{tabularx}{\textwidth}{|X|X|X|}
\hline
\textbf{Item} & \textbf{Total Reduction (Left)} & \textbf{Total Reduction (Right)} \\
\hline
Original & 50 + 100 / 2 = 100 & 75 + 50 / 2 = 100 \\
Simplified & 100 & 100 \\
\hline
\end{tabularx}
\end{table}

We also calculate the total degrees of myopia and astigmatism:

\begin{table}[h!]
\centering
\begin{tabularx}{\textwidth}{|X|X|X|}
\hline
\textbf{Item} & \textbf{Total (Left)} & \textbf{Total (Right)} \\
\hline
Original & 300 + 125 / 2 = 362.5 & 500 + 125 / 2 = 562.5 \\
Simplified & 360 & 560 \\
\hline
\end{tabularx}
\end{table}

An interesting observation arises when comparing the reduction in astigmatism to myopia. For example, when the left eye's myopia is reduced by 50 degrees and the astigmatism is reduced by 100 degrees, this is effectively equivalent to a 100-degree reduction in total refractive error. The right eye shows a similar pattern, where a 75-degree reduction in myopia, plus half the reduction in astigmatism (i.e., 50/2), adds up to an overall 100-degree reduction in refractive error. This approach underscores the interconnectedness of these two refractive conditions, with astigmatism reduction playing a pivotal role in the overall improvement in vision.

Furthermore, it is hypothesized that the amount of astigmatism reduction may be influenced by the initial severity of astigmatism. Those with more pronounced astigmatism may experience more substantial reductions due to the process of eyeball reshaping. This could reflect the timing and nature of the astigmatism's development, which is linked to the point in time when the original deformation occurred and began to correct itself.

It is also worth noting that the specific shape of each person's eye can affect the rate and regularity of this process. The progression of myopia and astigmatism may follow a regular pattern in some individuals, leading to a more uniform form of myopia, while in others, irregularities in the eye shape may lead to the development of astigmatism. Astigmatism can be conceptualized as a columnar lens, similar to the side of a wine bottle, where the deformation creates a cylindrical shape rather than a spherical one.

\section{The ``Just Barely Clear'' Principle}

\subsection{Personal Experience}
I would like to share a significant personal experience that underscores the importance of the ``just barely clear'' principle in the process of myopia reversal. During the first year of implementing this method, my myopia was reduced by approximately 100 degrees, which was a promising outcome. However, in the subsequent year, I observed little to no further improvement. Upon reflecting on this stagnation, I was able to identify the likely cause:

In the first year, my primary activity involved computer work, during which I wore glasses with a 150-degree reduction. This setup allowed me to frequently experience the ``just barely clear'' state, where vision was sufficiently clear without being overly sharp. However, in the second year, as I shifted my focus to studying for my associate degree at home, I reduced my computer usage and primarily relied on my mobile phone for reading materials and exercises.

This change in my daily habits led me to revisit the work of Yin Wang, particularly the passage quoted in Section~\ref{sec:background}, which proved to be a turning point in my understanding. Based on this understanding, I experimented with glasses offering 175- and 200-degree reductions, which provided a greater opportunity for the ``just barely clear'' state. I am hopeful that this adjustment will lead to a further reduction in my myopia by approximately 100 degrees over the course of the next year.

From this experience, I have learned that the key to effective myopia reversal lies not in adhering to a fixed degree of reduction, but in consistently maintaining the ``just barely clear'' state. The specific degree of reduction required to achieve this state will vary for each individual, but the critical factor is identifying and sustaining this optimal point of clarity, which facilitates the gradual adaptation necessary for myopia reduction.

\subsection{Defining the State}
The ``just barely clear'' state can be characterized as follows:
\begin{itemize}
    \item Objects should be sufficiently clear to view without straining.
    \item Text should be readable, but not completely crisp.
    \item The eyes should feel relaxed, not engaged in heavy focusing efforts.
\end{itemize}

\subsection{Prescription Reductions by Activity}
Through experimentation, the following general rules were found:
\begin{itemize}
    \item For mobile phone use, a reduction of 175-200 degrees achieved the ``just barely clear'' state.
    \item For computer use, a reduction of 150 degrees typically suffices.
    \item The ``just barely clear'' state encourages gradual improvement without forcing excessive clarity.
\end{itemize}

The optimal reduction varies between individuals, but the key takeaway is that forcing clarity too early or with too large a reduction can hinder natural adaptation and improvement. Moving to stronger reductions too quickly can result in stagnation, which led to the realization that the ``just barely clear'' state is crucial for continued myopia reversal.

\section{Practical Tips}
For those interested in trying the natural vision restoration method, here are some practical tips based on my experience:
\begin{itemize}
    \item \textbf{Initial Discomfort:} When you first switch to glasses with a reduced prescription (e.g., 150 degrees less), you may experience some discomfort, especially in terms of seeing distant objects clearly. However, the eyes adapt quickly, and the discomfort typically subsides within a few days.
    \item \textbf{Adjusting to Reduced Prescription:} Over time, wearing glasses with a reduced prescription for near vision becomes more comfortable and should not interfere with daily activities, such as working, studying, or using mobile devices.
    \item \textbf{Eyeglasses and Activities:} For activities such as driving, attending courses, or watching movies, you may want to consider adjusting your glasses prescription to be slightly higher. However, for everyday tasks such as using a computer or mobile phone, the reduced prescription should be sufficient.
    \item \textbf{Sustain the Optimal Point:} Rather than focusing on a fixed number of degrees of reduction, consistently maintain the personalized ``just barely clear'' threshold of clarity that encourages the eye to adapt without straining.
\end{itemize}

\section{Discussion}
The concept of ``just barely clear'' vision has proven crucial in natural vision restoration. Unlike the conventional approach of prescribed reductions based on specific degree reductions, this approach emphasizes gradual adaptation. The reduction does not need to follow a fixed number of degrees, but instead focuses on finding the threshold where objects are visible without excessive strain.

This concept aligns with the theory that myopia is the result of optical deformation, and by encouraging the eyes to work slightly harder to focus, the optical axis is gradually shortened, reversing myopia. The experience underscores the importance of balance in vision clarity---too much clarity results in less stimulation for the eyes to adapt, while too little clarity can lead to strain.

The underlying reason behind the success of this method remains complex and is still under research. According to Todd Becker's article, reducing the prescription of glasses encourages the eye muscles to adapt to a natural, relaxed state. While I do not fully understand the scientific mechanisms, the observed improvements in myopia and astigmatism suggest that this method is effective in some cases.

Interestingly, my experiment revealed that the recovery rate of my left and right eyes differed. It seems that my right eye, which had a higher degree of myopia, showed faster recovery, potentially due to more significant distortion in the eye muscles. This recovery pattern aligns with concepts in machine learning, where a system (in this case, the eye muscles) improves rapidly at first and then levels off over time. Similarly, for individuals with varying severities of myopia and astigmatism, the reduction process may take different forms: more severe myopia may see faster progress, while less pronounced myopia may require more time to yield visible results.

\section{Future Outlook}
I will continue monitoring my eyesight in the coming years and provide further reports. Based on the current trajectory, I expect my left eye to recover fully within three years, while my right eye may also continue to improve. I am also hopeful that the adjustments motivated by the ``just barely clear'' principle will lead to a further reduction in my myopia by approximately 100 degrees over the course of the next year.

\section{Conclusion}
The results of this experiment suggest that the natural vision restoration method is a promising approach for reducing myopia and astigmatism. This investigation into the relationship between astigmatism and myopia during the process of eyeball deformation and reversal supports the hypothesis that astigmatism reduction is linked to the physical process of myopia correction, with both conditions improving concurrently. Furthermore, the ``just barely clear'' principle highlights the importance of finding a personalized threshold of clarity that encourages the eye to adapt without straining; gradual, consistent reduction in prescription, adjusted according to individual needs, is more effective than adhering to a one-size-fits-all approach. Although further research is needed to fully understand the mechanisms behind this method, the results so far are encouraging and open the door for further investigation into individualized myopia reversal techniques.

\section*{References}
\begin{itemize}
    \item Becker, Todd. (2014). ``Myopia: A Modern Yet Reversible Disease.'' \url{https://gettingstronger.org/2014/08/myopia-a-modern-yet-reversible-disease}.
    \item Wang, Yin. (2022). ``Natural Vision Restoration Method.'' \url{https://www.yinwang.org/blog-cn/2022/02/22/myopia}.
\end{itemize}

\section*{Citation}
Li, Zhiwei. (2026 August). \textit{Natural Vision Restoration: Experimental Verification, Astigmatism Reduction, and the 'Just Barely Clear' Principle}. Zhiwei's Blog. \url{https://lzwjava.com/natural-vision-restoration-en}

Or in BibTeX:
\begin{verbatim}
@article{li2026naturalvisionrestoration,
  title   = "Natural Vision Restoration: Experimental Verification,
             Astigmatism Reduction, and the 'Just Barely Clear' Principle",
  author  = "Li, Zhiwei",
  journal = "lzwjava.com",
  year    = "2026",
  month   = "Aug",
  url     = "https://lzwjava.com/natural-vision-restoration-en"
}
\end{verbatim}

\end{document}
